%****************************** Preámbulo ******************************%
\documentclass[letterpaper, 12pt]{article}

% <<< BEGIN idioma y codificación

% Idioma de trabajo: español, "table" -> "tabla" y punto decimal en fórmulas (sin la opción sería coma).
\usepackage[spanish, es-tabla, es-nodecimaldot]{babel}

% Acentos, ñ y separación silábica correcta de palabras acentuadas.
\usepackage[utf8]{inputenc}
\usepackage[T1]{fontenc}
\usepackage{lmodern}

% >>> END

% <<< BEGIN paquetes y macros

% Imágenes y ubicación fija de figuras con \begin{figure}[H].
\usepackage{graphicx}
\usepackage{float}
\graphicspath{{images/}}

% Notación (derivadas, vectores, brakets, unidades) y entornos matemáticos.
\usepackage{physics, amsmath, amssymb}

% Tablas con líneas horizontales limpias (\toprule, \midrule, \bottomrule).
\usepackage{booktabs}

% Gráficos con el lenguaje nativo de LaTeX.
\usepackage{tikz}
\usetikzlibrary{fit, backgrounds, arrows.meta, positioning, babel}

% Colores para gráficas y diagramas TikZ (los mismos del template de Beamer).
% Los colores UNAM (UNAMDarkBlue, UNAMLightBlue, UNAMGold) vienen de unampaper.sty.
\definecolor{myblue}{RGB}{30,100,180}
\definecolor{myred}{RGB}{180,30,30}
\definecolor{mygreen}{RGB}{20,140,60}
\definecolor{mypurple}{RGB}{120,30,160}
\definecolor{myorange}{RGB}{220,120,0}

% Listas con formato configurable (mejor que enumerate a secas).
\usepackage{enumitem}

% Texto a dos columnas: \begin{multicols}{2} ... \end{multicols}.
\usepackage{multicol}
\setlength\columnsep{20pt}

% Fórmulas más grandes dentro del texto (útil en ejercicios resueltos).
\usepackage{relsize}
\newcommand{\mlr}[1]{\mbox{$\mathlarger{#1}$}}

% Unidades: 9.8\un{m/s$^2$}.
\newcommand{\un}[1]{\mbox{ #1}}

% Citas como superíndice.
\newcommand{\upcite}[1]{\mbox{$^\text{\cite{#1}}$}}

% >>> END

% <<< BEGIN estilo UNAM

% Márgenes, título, índice, PromptBlock y enlaces (colores UNAM solo en enlaces).
% Va al final del preámbulo porque carga hyperref.
\usepackage{unampaper}

% >>> END

% <<< BEGIN metadatos

\title{Título del trabajo}
\author{Nombre del autor}
% \date{Noviembre 2026}   % opcional

% >>> END

%****************************** Documento ******************************%
\begin{document}

    % Numeración romana desde la portada.
    \pagenumbering{Roman}

    \maketitle

    \begin{abstract}
        Resumen del trabajo: problema, método y resultado principal en un solo párrafo.
    \end{abstract}

    \unamindice

    \section{Teoría}
        \subsection{Subsección}
            Texto con una ecuación numerada y etiquetada:

            \begin{equation}\label{eq:ejemplo}
                \delta(\vb{x},t) = \frac{\rho(\vb{x},t) - \bar{\rho}(t)}{\bar{\rho}(t)}.
            \end{equation}

            Referencia a la ecuación \ref{eq:ejemplo} y cita en superíndice\upcite{ref1}. Ecuaciones alineadas:

            \begin{align}
                \delta(\vb{k}) &= \sum_{n=1}^{\infty} \delta^{(n)}(\vb{k}), \\
                \expval{\delta(\vb{k})\delta(\vb{k}')} &= (2\pi)^3 \delta_D(\vb{k}+\vb{k}') P(k).
            \end{align}

        % Con matemáticas en el título usar \texorpdfstring{con LaTeX}{texto plano para el PDF}.
        \subsection{\texorpdfstring{Subsección con $F_2$ en el título}{Subsección con F2 en el título}}
            Texto de la subsección con una cita normal \cite{ref2}.

    \section{Cómputo}
        \subsection{Método numérico}
            Texto con \texttt{código}, unidades (9.8\un{m/s$^2$}) y listas:

            \begin{itemize}
                \item Primer nivel.
                    \begin{itemize}
                        \item Segundo nivel.
                            \begin{itemize}
                                \item Tercer nivel.
                            \end{itemize}
                    \end{itemize}
                \item Otro elemento.
            \end{itemize}

            \begin{enumerate}
                \item Primer paso.
                \item Segundo paso.
            \end{enumerate}

        \subsection{Subsección}
            \subsubsection{Subsubsección}
                Texto de la subsubsección.

    \section{Resultados}
        Referencia a la Figura~\ref{fig:ejemplo} y a la Tabla~\ref{tab:ejemplo}.

        \begin{figure}[H]
            \centering
            % reemplazar por: \includegraphics[width=0.85\linewidth]{mi-figura.pdf}
            \begin{tikzpicture}
                \draw[thick, draw=black!60, fill=black!3] (0,0) rectangle (10,4);
                \node[text=black!60] at (5,2) {Figura};
            \end{tikzpicture}
            \caption{Descripción de la figura.}
            \label{fig:ejemplo}
        \end{figure}

        \begin{table}[H]
            \centering
            \caption{Descripción de la tabla.}
            \label{tab:ejemplo}
            \begin{tabular}{lcc}
                \toprule
                Parámetro & Valor & Error \\
                \midrule
                $\Omega_{m0}$ & 0.31 & 0.01 \\
                $h$           & 0.67 & 0.02 \\
                \bottomrule
            \end{tabular}
        \end{table}

    \section{Uso de los modelos LLM}
        Modelos utilizados, versiones y enlaces a los historiales:

        \begin{enumerate}
            \item \url{https://claude.ai/share/...}
        \end{enumerate}

        Instrucciones de sistema utilizadas:

        \begin{PromptBlock}
SYSTEM INSTRUCTIONS

Texto literal; respeta espacios y saltos de línea, y puede continuar
de una página a otra.
        \end{PromptBlock}

    \section{Conclusiones}
        Texto de las conclusiones.

    \begin{thebibliography}{0}

        \bibitem{ref1} Apellido, A., \& Apellido, B. (Año). Título del artículo. \textit{Revista}, volumen(número), página. \url{https://doi.org/...}

        \bibitem{ref2} Apellido, A. (Año). Título del preprint. arXiv preprint arXiv:XXXX.XXXXX. \url{https://arxiv.org/abs/XXXX.XXXXX}

    \end{thebibliography}

\end{document}
